\documentclass[tikz,border=6pt]{standalone}
\usepackage{ctex}
\usepackage{amsmath}
\usepackage{amssymb}
\usepackage{pgfplots}
\pgfplotsset{compat=1.18}
\usetikzlibrary{calc,angles,quotes,arrows.meta,positioning,decorations.markings}
\begin{document}
\begin{tikzpicture}
\begin{axis}[
    width=12.5cm, height=9cm,
    xlabel={$x$}, ylabel={$y$},
    xmin=-2, xmax=2.6, ymin=0, ymax=8,
    samples=200,
    axis lines=left,
    grid=major, grid style={gray!18},
    tick label style={font=\small},
    label style={font=\small},
    legend pos=north west, legend cell align=left,
]
% 函数 y=e^x
\addplot[blue, very thick, domain=-2:2.07] {exp(x)};
\addlegendentry{$y=e^x$}
% 切线：在 x0=1 处 y = e + e(x-1) = e*x
\addplot[red, very thick, domain=-0.2:2.4] {2.718281828*x};
\addlegendentry{切线 $y=e\,x$，斜率 $=e$}
% 切点 (1, e)
\addplot[only marks, mark=*, mark size=2pt, black] coordinates {(1,2.718281828)};
\node[anchor=south east, font=\small] at (axis cs:1,2.7183) {$P(1,\,e)$};
% 直角三角形：从 (1,0) 上到 (1,e) 再水平表示“高度=斜率”
% 竖直高度段：从 (1,0) 到 (1,e)
\draw[densely dashed, gray] (axis cs:1,0) -- (axis cs:1,2.718281828);
% 切线的斜率三角形：从 (1,e) 水平到 (2,e)，再竖直到 (2,2e)
\draw[red!70!black, thick] (axis cs:1,2.718281828) -- (axis cs:2,2.718281828);
\draw[red!70!black, thick] (axis cs:2,2.718281828) -- (axis cs:2,5.436563656);
% 直角标记
\draw[red!70!black] (axis cs:1.85,2.718281828) -- (axis cs:1.85,2.868281828) -- (axis cs:2,2.868281828);
% 标注 水平=1, 竖直=e
\node[red!70!black, anchor=north, font=\small] at (axis cs:1.5,2.69) {$\Delta x=1$};
\node[red!70!black, anchor=west, font=\small] at (axis cs:2.02,4.1) {$\Delta y=e$};
% 高度标注
\node[blue!60!black, anchor=east, font=\small] at (axis cs:0.97,1.5) {高度 $=e^1=e$};
% 核心注释框
\node[draw=red!55, fill=red!8, fill opacity=0.85, text opacity=1,
      rounded corners, align=left, font=\small, inner sep=4pt, anchor=south west]
  at (axis cs:-1.95,3.05)
  {在 $P(1,e)$ 处：\\
   切线斜率 $=\dfrac{\Delta y}{\Delta x}=\dfrac{e}{1}=e$\\[2pt]
   恰等于该点高度 $e^1=e$\\
   即 $(e^x)'=e^x$};
\end{axis}
\end{tikzpicture}
\end{document}
